\metatitle{Plane partitions}
\metadescription{Plane partitions.}


\section[planePartitions]{Plane partitions}

A \defin{plane partition} is an array $\pi_{ij}$ of non-negative 
integers with weakly decreasing rows and columns.
We let $|\pi|$ denote the total sum of the entries.

See \href{http://www-math.mit.edu/~rstan/pubs/pubfiles/75.pdf}{Stanley's survey}
for a bit more background.

\begin{theorem}[MacMahon, \cite{MacMahon1960}]
We have
\[
\sum_{\pi} q^{|\pi|} = 
\prod_{i=1}^a \prod_{j=1}^b \prod_{k=1}^m
\frac{1-q^{i+j+k-1}}{1-q^{i+j+k-2}}
\]
where the sum is over all plane partitions 
with $\leq a$ rows, $\leq b$ columns and largest entry $\leq m$.
\end{theorem}

\begin{example}
For \(a=1\), \(b=2\), and \(m=2\), the plane partitions are weakly
decreasing rows \((u,v)\) with \(2\geq u\geq v\geq 0\).  They are
\[
  (0,0),\quad (1,0),\quad (1,1),\quad
  (2,0),\quad (2,1),\quad (2,2),
\]
so the generating polynomial is
\[
  1+q+2q^2+q^3+q^4.
\]
The product formula gives the same value:
\[
  \frac{1-q^2}{1-q}
  \frac{1-q^3}{1-q^2}
  \frac{1-q^3}{1-q^2}
  \frac{1-q^4}{1-q^3}
  =
  \frac{(1-q^3)(1-q^4)}{(1-q)(1-q^2)}
  =
  1+q+2q^2+q^3+q^4.
\]
\end{example}


\begin{theorem}[Kreweras, \cite{Kreweras1965}]
The number of plane partitions of shape $\lambda/\mu$
with maximal entry at most $m$, is given by
\[
\det\left(
\binom{\lambda_i-\mu_j + m}{i-j+m}
\right)_{i,j=1,\dotsc,\ell}
\]
where $\ell = \length(\lambda)$.
\end{theorem}
Note that this allows us to efficiently compute the
\hyperref[ehrhartPolynomialDefinition]{Ehrhart polynomial} of
\hyperref[orderPolytopeDefinition]{order polytopes} associated with Ferrers
diagrams.

\name{Ilse Fischer} and \name{Florian Schreier-Aigner} introduce a family of
Schur-positive symmetric functions defined as sums over totally symmetric
plane partitions \cite{FischerSchreierAigner2024ASM}.  For one specialization
this family agrees with a multivariate generating function for extended
alternating sign matrices.
\name{Sam Hopkins} and \name{Tri Lai} prove a product formula for plane
partitions of shifted double staircase shape \cite{HopkinsLai2021}.
\name{Igor Pak} and \name{Fedor Petrov} study hidden symmetries of weighted
lozenge tilings \cite{PakPetrov2020}.

\name{Sam Hopkins}: \href{https://mathoverflow.net/questions/350445/determinantal-formula-for-plane-partitions-of-shifted-shape}{Is there a similar formula for shifted shapes?}
